% 量子物理基础填空题12.2.1—12.2.13；用户原书/PDF第317—318页。
\subsection{填空题：光电截止、氢谱与测不准关系}

\sourceproblem{12.2.1}{317}{317}
\problemtext 波长300 nm照射某金属，光电子动能由0到$4.0\times10^{-19}$ J，
求遏止电压及红限频率。
\workings
由$eU_a=K_{\max}$，有
\[
\boxed{U_a=\frac{4.0\times10^{-19}}{1.602176634\times10^{-19}}
\approx2.50\,\mathrm V}.
\]
光子能$hc/\lambda\approx6.62\times10^{-19}$ J，
逸出功$W=hc/\lambda-K_{\max}$，
\[
\boxed{\nu_0=W/h
=c/\lambda-K_{\max}/h\approx3.96\times10^{14}\,\mathrm{Hz}}.
\]
\crosscheck 红限频率必须为正且低于入射光频率$10^{15}$ Hz。

\sourceproblem{12.2.2}{317}{317}
\problemtext X射线分别在$45^\circ$及$60^\circ$发生康普顿散射，
比较散射波长增量$\Delta\lambda_1:\Delta\lambda_2$。
\workings
由$\Delta\lambda=\lambda_C(1-\cos\theta)$，
\[
\boxed{\frac{\Delta\lambda_1}{\Delta\lambda_2}
=\frac{1-\cos45^\circ}{1-\cos60^\circ}
=2-\sqrt2\approx0.5858}.
\]
\crosscheck 60度散射增量更大，所得比值小于1。

\sourceproblem{12.2.3}{317--318}{317--318}
\problemtext 氢原子电子由能量$-0.85$ eV跃迁到$-3.4$ eV；
求释放光子能量及两状态主量子数。
\workings
$E_n=-13.6/n^2$ eV，$-0.85$ eV对应$n=4$，
$-3.4$ eV对应$n=2$，
\[
\boxed{E_\gamma=3.4-0.85=2.55\,\mathrm{eV}},
\qquad \boxed{n=4\to2}.
\]
\crosscheck 波长约$1240/2.55=486$ nm，是巴耳末可见谱线。

\sourceproblem{12.2.4}{318}{318}
\problemtext 依据原图$n=1,2,3,4$的能级及箭头，
选出所发光子波长最短、频率最低的跃迁。
\workings
波长最短对应能级差最大，故
$\boxed{n=4\to n=1}$；
频率最低对应所画向下跃迁中能级差最小，故
$\boxed{n=4\to n=3}$。
\crosscheck 这里比较的是图中\emph{发射箭头}，并非任意能级差或吸收跃迁。

\sourceproblem{12.2.5}{318}{318}
\problemtext 氢原子基态能量$E_1=-13.6$ eV，求$n=4$能量。
\workings
\[
\boxed{E_4=-\frac{13.6}{4^2}=-0.85\,\mathrm{eV}}.
\]
\crosscheck 相对电离极限零能，束缚态能量均为负。

\sourceproblem{12.2.6}{318}{318}
\problemtext 某电子在原子激发态中平均停留$10^{-8}$s；
估算能级能量最小不确定度。
\workings
按标准能时估计$\Delta E\,\Delta t\gtrsim\hbar/2$，
\[
\boxed{\Delta E\gtrsim\frac{\hbar}{2(10^{-8})}
=5.27\times10^{-27}\,\mathrm J
\approx3.29\times10^{-8}\,\mathrm{eV}}.
\]
\crosscheck 能时不确定度不同于位置--动量算符的严格同类形式，
自然线宽的定义也可能采用$\Gamma\simeq\hbar/\tau$
而差一个约为2的系数；应说明采用的约定。

\sourceproblem{12.2.7}{318}{318}
\problemtext 用波长$\lambda$表达光子的能量、动量大小与相对论质量等效值。
\workings
\[
\boxed{\varepsilon=\frac{hc}{\lambda}},\qquad
\boxed{p=\frac h\lambda},\qquad
\boxed{m_{\rm eq}=\frac{\varepsilon}{c^2}=\frac h{\lambda c}}.
\]
\crosscheck 光子的\emph{静质量为零}；$m_{\rm eq}$仅为原书旧术语下的能量等效质量。

\sourceproblem{12.2.8}{318}{318}
\problemtext 氢原子电子主量子数$n=3$，求允许轨道角动量大小最大值。
\workings
$l=0,1,2$，最大$l=n-1=2$，
\[
\boxed{L_{\max}=\sqrt{l(l+1)}\hbar=\sqrt6\,\hbar}.
\]
\crosscheck 不用玻尔旧量子论的$L=n\hbar$取代现代量子角动量。

\sourceproblem{12.2.9}{318}{318}
\problemtext 电子限于宽度$\Delta x=5$ nm的区域，
按\emph{原题指定的估算关系}$\Delta x\Delta p_x\ge h$，
求动量分量的最小估计不确定度。
\workings
依题面特定估算式
\[
\boxed{\Delta p_x\gtrsim\frac h{\Delta x}
=\frac{6.62607015\times10^{-34}}{5\times10^{-9}}
\approx1.33\times10^{-25}\,\mathrm{kg\,m/s}}.
\]
\crosscheck \textbf{源式与标准严格界限不同：}标准海森堡不等式为
$\Delta x\Delta p_x\ge\hbar/2=h/(4\pi)$。
若按严格标准只得$1.06\times10^{-26}$ kg\,m/s，
本条为忠于\emph{原题指定公式}而计算，不能混称两式等价。

\sourceproblem{12.2.10}{318}{318}
\problemtext 用波函数$\Psi(\bm r,t)$描述微观粒子，
写出$\Psi\Psi^*$的概率解释、波函数要求及归一化条件。
\workings
\[
\boxed{\rho(\bm r,t)=\Psi\Psi^*=|\Psi|^2},\qquad
\boxed{\int_{\mathbb R^3}|\Psi(\bm r,t)|^2\,d^3r=1}.
\]
波函数需单值、有限、满足问题的连续/可微与边界条件。
\crosscheck 若只描述一维粒子，则归一化体积分应改为$\int|\psi(x)|^2dx$。

\sourceproblem{12.2.11}{318}{318}
\problemtext 电子自旋磁量子数$m_s$可能的两个值。
\workings
\[
\boxed{m_s=+\frac12}\qquad\boxed{m_s=-\frac12}.
\]
\crosscheck 电子自旋量子数为$s=1/2$，而$m_s$为其投影量子数。

\sourceproblem{12.2.12}{318}{318}
\problemtext 主量子数$n=2$的一个电子壳层最多有多少种单电子量子态？
\workings
计入自旋的态数为$2n^2$，
\[
\boxed{N=2(2)^2=8}.
\]
\crosscheck $n=2$由$2s$（2个态）和$2p$（6个态）组成。

\sourceproblem{12.2.13}{318}{318}
\problemtext 多电子原子中电子排布遵从哪两项基本原则？
\workings
依原书填\textbf{泡利不相容原理}与\textbf{能量最小原理}。
\crosscheck 实际多电子组态还需结合洪特规则等判断简并亚层占据方式；
本题只要求两项教材概括原则。
